% !TeX program = xelatex
% ============================================================
% AI-assisted C/C++ development workflow
% XeLaTeX + TikZ
% Design goals:
%   1. Styles and content are separated as much as possible.
%   2. Nodes use relative positioning (positioning library), no absolute coords.
%   3. Most later edits only need changes in SECTION 3 (CONTENT).
% ============================================================

\documentclass[11pt,a4paper]{article}

% -------------------- 1. PACKAGES ---------------------------
\usepackage[a4paper,margin=14mm]{geometry}
\usepackage{fontspec}
\usepackage{xeCJK}
\usepackage{xcolor}
\usepackage{tikz}
\usetikzlibrary{
  arrows.meta,
  positioning,
  shapes.geometric,
  fit,
  backgrounds,
  calc
}
\usepackage{enumitem}
\usepackage{microtype}
\pagestyle{empty}

% -------------------- 2. TYPOGRAPHY + STYLE -----------------
% Replace these two lines if you prefer another installed Chinese font.
\setmainfont{DejaVu Sans}
\setCJKmainfont{Noto Sans CJK SC}
\setsansfont{DejaVu Sans}
\setCJKsansfont{Noto Sans CJK SC}

% ---- Palette: change colors here, not inside drawing code ----
\definecolor{HumanBlue}{HTML}{2563EB}
\definecolor{AIPurple}{HTML}{7C3AED}
\definecolor{ToolGreen}{HTML}{059669}
\definecolor{ReviewOrange}{HTML}{EA580C}
\definecolor{PassGreen}{HTML}{15803D}
\definecolor{FailRed}{HTML}{DC2626}
\definecolor{Ink}{HTML}{1F2937}
\definecolor{SoftGray}{HTML}{F3F4F6}
\definecolor{LineGray}{HTML}{64748B}
\definecolor{PaleBlue}{HTML}{EFF6FF}
\definecolor{PalePurple}{HTML}{F5F3FF}
\definecolor{PaleGreen}{HTML}{ECFDF5}
\definecolor{PaleOrange}{HTML}{FFF7ED}

% ---- Global TikZ knobs ----
\tikzset{
  >={Latex[length=2.8mm,width=1.9mm]},
  every node/.style={font=\sffamily},
  workflow/.style={
    draw=Ink!55,
    line width=0.85pt,
    rounded corners=2.5mm,
    text=Ink,
    align=center,
    minimum height=14mm,
    text width=72mm,
    inner xsep=5mm,
    inner ysep=3mm,
  },
  human/.style={workflow, draw=HumanBlue, fill=PaleBlue},
  ai/.style={workflow, draw=AIPurple, fill=PalePurple},
  tool/.style={workflow, draw=ToolGreen, fill=PaleGreen},
  review/.style={workflow, draw=ReviewOrange, fill=PaleOrange},
  decision/.style={
    diamond,
    aspect=2.25,
    draw=ReviewOrange,
    fill=PaleOrange,
    line width=0.9pt,
    align=center,
    text=Ink,
    inner sep=1.5mm,
    text width=27mm,
  },
  terminal/.style={
    rounded corners=7mm,
    draw=PassGreen,
    fill=PassGreen!8,
    line width=1pt,
    minimum width=40mm,
    minimum height=11mm,
    align=center,
    text=Ink,
  },
  flow/.style={->, draw=LineGray, line width=1.05pt},
  feedback/.style={->, draw=FailRed, line width=1.05pt, dashed},
  passflow/.style={->, draw=PassGreen, line width=1.05pt},
  rolebox/.style={
    rounded corners=2mm,
    draw=Ink!25,
    fill=white,
    line width=0.6pt,
    inner sep=3mm,
    text width=48mm,
    align=left,
    font=\sffamily\small,
  },
  phasebg/.style={
    rounded corners=3mm,
    draw=none,
    fill=SoftGray,
    inner sep=4mm,
  },
  tag/.style={
    rounded corners=1.5mm,
    inner xsep=2.2mm,
    inner ysep=1.0mm,
    font=\sffamily\scriptsize\bfseries,
    text=white,
  }
}

% -------------------- 3. CONTENT ----------------------------
% Most future maintenance should happen in this section only.

\newcommand{\DocTitle}{AI 辅助 C/C++ 开发工作流}
\newcommand{\DocSubtitle}{人定义约束，AI 高效执行，工具提供证据，人负责最终判断}

\newcommand{\HumanDefineTitle}{人类工程师：定义问题}
\newcommand{\HumanDefineBody}{需求、接口、性能目标、生命周期、并发约束与系统不变量}

\newcommand{\AIDesignTitle}{AI：参与设计讨论}
\newcommand{\AIDesignBody}{补充方案、暴露遗漏、比较实现路径；架构取舍仍由人主导}

\newcommand{\AICodeTitle}{AI：生成与重构代码}
\newcommand{\AICodeBody}{脚手架、初版实现、重复性代码、测试框架、重构候选方案}

\newcommand{\StaticTitle}{编译器与静态检查}
\newcommand{\StaticBody}{Warnings / clang-tidy / 格式检查 / 编译期约束}

\newcommand{\RuntimeTitle}{运行期验证}
\newcommand{\RuntimeBody}{ASan + UBSan；并发场景单独使用 TSan}

\newcommand{\TestTitle}{测试与测量}
\newcommand{\TestBody}{单元测试、边界测试、Stress / Fuzz、Benchmark / perf}

\newcommand{\ReviewTitle}{人工 Review}
\newcommand{\ReviewBody}{检查所有权、生命周期、锁粒度、失败路径、实时性与性能假设}

\newcommand{\DecisionText}{是否满足\\约束与证据？}
\newcommand{\PassText}{通过：合入 / 发布}
\newcommand{\FailText}{问题与证据反馈给 AI\\修改后重新验证}

\newcommand{\HumanRole}{%
  \textbf{人负责：}\\[-1mm]
  \begin{itemize}[leftmargin=*,nosep]
    \item 定义 Contract / Invariant
    \item 架构与接口取舍
    \item 生命周期、并发与性能判断
    \item 最终是否可合入
  \end{itemize}}

\newcommand{\AIRole}{%
  \textbf{AI 负责：}\\[-1mm]
  \begin{itemize}[leftmargin=*,nosep]
    \item Boilerplate 与初版实现
    \item 设计讨论与重构建议
    \item 测试初稿与日志分析
    \item 编译错误与工具输出解释
  \end{itemize}}

\newcommand{\ToolRole}{%
  \textbf{工具负责：}\\[-1mm]
  \begin{itemize}[leftmargin=*,nosep]
    \item 提供可复现证据
    \item 提高错误发现率
    \item 验证正确性与性能假设
  \end{itemize}}

% -------------------- 4. DRAWING ----------------------------
\begin{document}

\begin{center}
  {\LARGE\bfseries\color{Ink}\DocTitle}\par
  \vspace{1.5mm}
  {\small\color{LineGray}\DocSubtitle}
\end{center}

\vspace{2mm}

\begin{tikzpicture}[node distance=6.5mm and 10mm]

  % ---- Main vertical flow: all positions are relative ----
  \node[human] (define) {\textbf{\HumanDefineTitle}\\[1mm]\small \HumanDefineBody};

  \node[ai, below=of define] (design) {\textbf{\AIDesignTitle}\\[1mm]\small \AIDesignBody};

  \node[ai, below=of design] (code) {\textbf{\AICodeTitle}\\[1mm]\small \AICodeBody};

  \node[tool, below=of code] (static) {\textbf{\StaticTitle}\\[1mm]\small \StaticBody};

  \node[tool, below=of static] (runtime) {\textbf{\RuntimeTitle}\\[1mm]\small \RuntimeBody};

  \node[tool, below=of runtime] (tests) {\textbf{\TestTitle}\\[1mm]\small \TestBody};

  \node[review, below=of tests] (review) {\textbf{\ReviewTitle}\\[1mm]\small \ReviewBody};

  \node[decision, below=8mm of review] (decision) {\textbf{\DecisionText}};

  \node[terminal, below=8mm of decision] (pass) {\textbf{\PassText}};

  % ---- Straight flow arrows ----
  \draw[flow] (define) -- (design);
  \draw[flow] (design) -- (code);
  \draw[flow] (code) -- (static);
  \draw[flow] (static) -- (runtime);
  \draw[flow] (runtime) -- (tests);
  \draw[flow] (tests) -- (review);
  \draw[flow] (review) -- (decision);
  \draw[passflow] (decision) -- node[right, font=\scriptsize\bfseries, text=PassGreen] {是} (pass);

  % ---- Feedback loop: relative routing, no absolute coordinates ----
  % Keep the feedback card on the right, but route the return path BELOW the
  % main flow and then UP the left side.  This deliberately avoids crossing
  % the AI/TOOLS responsibility cards on the right.
  \node[review, text width=49mm, right=18mm of decision] (feedback) {\textbf{发现问题}\\[1mm]\small \FailText};
  \draw[feedback] (decision.east) -- node[above, font=\scriptsize\bfseries, text=FailRed] {否} (feedback.west);

  % Helper coordinates are all derived from node anchors, so maintenance
  % remains relative-position based rather than relying on absolute (x,y).
  \coordinate (feedbackbottom) at ([yshift=-7mm]pass.south);
  \coordinate (feedbackleft)   at ([xshift=-9mm]code.west);
  \draw[feedback]
    (feedback.south)
    |- (feedbackbottom)
    -| (feedbackleft)
    -- node[above, near end, font=\scriptsize\bfseries, text=FailRed] {返回修改} (code.west);

  % ---- Role summary panel, also relatively positioned ----
  \node[rolebox, right=18mm of design] (humanrole) {\HumanRole};
  \node[rolebox, below=5mm of humanrole] (airole) {\AIRole};
  \node[rolebox, below=5mm of airole] (toolrole) {\ToolRole};

  % ---- Small role tags ----
  \node[tag, fill=HumanBlue, above left=-1mm and 0mm of humanrole.north west] {HUMAN};
  \node[tag, fill=AIPurple, above left=-1mm and 0mm of airole.north west] {AI};
  \node[tag, fill=ToolGreen, above left=-1mm and 0mm of toolrole.north west] {TOOLS};

  % ---- Soft background grouping for main stages ----
  \begin{scope}[on background layer]
    \node[phasebg, fit=(define)] {};
    \node[phasebg, fit=(design)(code)] {};
    \node[phasebg, fit=(static)(runtime)(tests)] {};
    \node[phasebg, fit=(review)(decision)(pass)(feedback)] {};
  \end{scope}

\end{tikzpicture}

\vfill
\begin{center}
  \small\bfseries\color{Ink}
  人定规，AI 施工；工具验收，人签字。
\end{center}

\end{document}
